\begin{enumerate}[label=\textbf{\ansnum.\arabic*}]
  \item $R=0.0175\cdot8/A$. $2.5$\,mm$^2$: $R=0.056\,\Omega$, $\Delta U=1.68$\,V ($3.5\%$). $4$\,mm$^2$: $0.035\,\Omega$, $1.05$\,V ($2.2\%$). $6$\,mm$^2$: $0.0233\,\Omega$, $0.70$\,V ($1.5\%$). The smallest with less than $3\%$ ($1.44$\,V) is $4$\,mm$^2$.
  \item $P=I^2R=900\cdot0.0233\approx21$\,W.
  \item $E=mgh=20\cdot9.81\cdot0.3\approx59$\,J. One tenth is $5.9$\,J in $0.2$\,s, about $29$\,W on average (with much higher short peaks). A battery takes the energy like charging; a power supply that cannot sink current lets the voltage rise on its output capacitors until its protection switches off, or the robot's electronics see over-voltage.
  \item Per session the battery would need about $3$ charges (one change per hour, the first battery is in already: $2$ changes). Two sessions: about $4$ changes and $6$ full charges ($6\cdot240=1440$\,Wh) per week, plus the waiting time, are saved.
  \item The fuse protects the \emph{cable}: if the cable is damaged anywhere along its length (pinched, cut), the short circuit is between the fuse and the damage only if the fuse sits at the source. A fuse at the robot end leaves the whole cable unprotected.
\end{enumerate}
